\section{FPGA and Synthesis Notes}

\subsection{Why simulation code and FPGA code differ}

A simulation memory can return data combinationally. FPGA block RAM usually returns data one clock after the address is registered. That difference strongly affects CPU architecture.

This is why the repository has both single-cycle and multi-cycle versions. The single-cycle core is easy to teach. The multi-cycle core is more realistic for FPGA memories.

\subsection{FPGA top modules}

The FPGA top-level files connect the SoC to board pins and clock/reset resources:

\begin{lstlisting}[language={}]
fpga_top.v
fpga_top_full.v
fpga_top_rtos.v
fpga_top_uart.v
\end{lstlisting}

The constraints directory contains Zynq-related \texttt{.xdc} files.

\subsection{Practical synthesis concerns}

The biggest synthesis concerns are:

\begin{itemize}
\item combinational divider cost in \texttt{alu.v};
\item combinational MMU walk in \texttt{mmu.v} if mapped directly;
\item memory timing assumptions in single-cycle designs;
\item timing closure in pipelined variants with forwarding and branch logic;
\item UART clock divisor correctness for the actual FPGA clock.
\end{itemize}
The multi-cycle MMU and multi-cycle CPU variants are better starting points for FPGA-friendly memory behavior.

\bigskip
